
\scanpage{55}
\mspage{11}
3) \deleted{Gommage} \inserted{Bouchage de trous, ou ``gommage''} de cercles de recollement.

Cette opération est relative à la donnée d'une 1-sous-variété fermée de $\partial X$, i.e. d'une partie $A'_l$ de $\pi_0(\partial X)\simeq A_l$ (ens des arêtes libres du graphe).
Elle consiste à boucher les trous correspondants, au moyen des cônes sur les composantes connexes du bord de $\Sigma'_l$.
Au niveau du graphe, cela consiste à gommer les arêtes libres correspondantes, avec bien entendu aussi les repères de marquage qui se trouvent au dessus --- \deleted{mais} les pondérations des sommets restent inchangées.

\deleted{\unsure{4}}
\uline{Gommage simultané mixte}, \deleted{\unsure{ou}} correspondant à une partie
\deleted{$A\amalg R'$} \inserted{$R'\amalg\Sigma'$} \deleted{quelconque} de $R\amalg\Sigma$,
qui est à la fois ouverte et fermée \inserted{et telle}
\deleted{\unsure{de} $A\amalg\vec R$},
\inserted{que $R'\supset R\cap\Sigma'$},
\deleted{ens. de composantes connexes de $R\amalg\Sigma$},
ou ce qui revient au $=$, une partie \inserted{$R'\amalg A'$} de $R\amalg A$, telle que $R'\supset R\mid A'$.

Au niveau des graphes pondérés à repères, cette opération consiste à exclure \deleted{les} \inserted{de $\vec R$} la partie corr. $\vec R'$, à contracter en un point les arêtes de $A'_{\deleted{\unsure{pq}}}$ (aussi bien libres que \unsure{liées}), en faisant suivre la pondération.
\mspage{54}
